\documentclass[a4paper,11pt]{ltjsarticle}
\usepackage{amsmath, amssymb, amsthm}
\usepackage[margin=1in]{geometry}
\usepackage{hyperref}

\newtheorem{theorem}{定理}[section]

\title{離散ミンコフスキー間隔におけるガウス整数を用いた正整数解の解法}
\author{
    \textbf{Noriaki Kihara (木原 範昭)} \\
    \textit{Chief Architect / CEO, WF System Co., Ltd.} \\
    \textit{Osaka University, School of Engineering Science (Alumni)} \\
    \textit{(Pioneer in Computational Geometry and 3D CAD Architecture)}
}
\date{}

\begin{document}
\maketitle

\section*{Acknowledgments}
本論文の著者は、計算ロジック構成と概念的パラダイムシフトの翻訳に貢献したAI技術「AntiGravity」に感謝の意を表する。

\begin{abstract}
本稿では、特殊相対性理論における平坦ミンコフスキー時空を、離散的な整数格子 $\mathbb{Z}^{1+N}$ 上で記述した場合に生じる数論的制約を考察する。空間的等方性を仮定すると、ミンコフスキー間隔の不変量はディオファントス方程式
\[ l^2 + t^2 = N x^2 \]
（$l, t, x \in \mathbb{Z}^+$）に帰着される。
この方程式の左辺は複素数平面上のガウス整数のノルムの構造を持ち、その正整数解の存在条件はフェルマーの二平方和定理として数学的に完全に解明されている。本論文の主眼は新たな数学的公式の発見ではなく、この\textbf{「既知のガウス整数の数論的制約」が、離散ミンコフスキー時空における「時空次元の選択」と見事にリンクしている点を見出したこと}にある。
証明の帰結として、$N$ の素因数分解において $p \equiv 3 \pmod{4}$ の素数の指数がすべて偶数であることが必要十分条件となるため、我々の住む4次元時空（$N=3$）では非自明解が存在しない（算術的アノマリー）ことが示される。
一方、$N=10$（$D=11$）や $N=25$（$D=26$）など、超弦理論やM理論で知られる特定の次元では解が存在する。本稿はこのような次元の一致を、連続理論における共形異常の消失と離散モデルにおける整数解の存在という観点から指摘するが、本モデルはあくまで特殊相対性理論の平坦極限における離散数論的制約を扱うものであり、量子重力理論の構築、ダイナミクス、ラグランジアン、または連続時空への直接的な拡張を意図するものではない。
本研究は、純粋数学における既知の算術的制約（ガウス整数の性質）が、離散物理モデルにおいて時空次元の選択に強い数学的制約を課すことを示し、平坦ミンコフスキー時空を整数格子上で離散化した場合における厳密解を与えるものである。
\end{abstract}

\section{導入 (Introduction)}
導きたい方程式は、
\[ l^2 = x^2 + y^2 + z^2 - t^2 \]
の \( l, x, y, z, t \) は反転対称正があると仮定しているので全て、0または正の整数解である。

さらに明記すべき点として、本稿は量子重力や弦理論などの本格的な物理理論を構築するものではない。あくまで、特殊相対性理論における平坦なミンコフスキー時空の極限を、離散的な整数格子 $\mathbb{Z}^{1+N}$ 上で記述した場合に現れるディオファントス方程式
\[ l^2 + t^2 = N x^2 \]
を考察するものであり、平坦ミンコフスキー時空を整数格子上で離散化した場合における厳密解を与えるものである。本方程式は曲率や相互作用を含まず、力学的なダイナミクスやラグランジアンは一切持たない。この制約された条件下において、等方的な測地線に対応する正整数解の存在条件を調べるのが本研究の主眼である。

\section{カードゲームのメタファー (Card Game Metaphor)}
ここでは、無限に広がったカードゲームのテーブルを考える。
まず、ディーラを \( l \) とする。

次に、次元をカードゲームのプレーヤとメタフォする。\\
0:次元　プレーヤは存在しない\\
1:次元　プレーヤは \( x \)\\
2:次元　プレーヤは \( x, y \)\\
3:次元　プレーヤは \( x, y, z \)

まず初期値（ゲームの開始時点）で、ディラーも、プレーヤも手札は0枚とする。

\subsection*{0:次元の盤面：}
プレーヤは存在しないので、ディーラー・テーブルの \( l^2 \) は２つの状態しかない。\\
カードが存在しない 0\\
カードが存在する  1

\subsection*{1:次元の盤面：}
プレーヤは \( x \) のみ。
\( x \) はディーラから任意の枚数のカードをもらえるが、
条件は、自分の手札は必ず平方数でなければならない。
また、ディーラーも、\( x \) の指定枚数と同じ枚数のカードを自身の手札に追加するが、ディーラーの手札も必ず平方数でなければならない。
このルールに違反すると、ゲームオーバーとなる。
この場合はごくシンプルで、
\[ l^2 = x^2 \]
となり、\( x \) の手札の枚数も、ディーラーの手札の枚数も、
0, 1, 4, 9 の平方数となる。また、\( x \) は手札が平方数となるカードをディーラーに要求できるので、ゲームは無限に終わらない。

\subsection*{2:次元の盤面：}
プレーヤは \( x, y \) となるが空間軸の対称性を仮定するので、特殊なルールがあり、\( y \) は、\( x \) と同じ数の手札しか要求できないとする。
つまり、
\[ l^2 = x^2 + y^2 = 2x^2 \]
となる。
\( 2x^2 \) が、平方数でなければゲームオーバになるので、\( x^2 \) は2の倍数でなければならないが、このルールになる枚数を \( x \) は必ずディーラに要求できるので、ゲームは無限に終わらない。

\subsection*{3:次元の盤面：}
プレーヤは、\( x, y, z \) となるが、同様に考えると、
\[ l^2 = x^2 + y^2 + z^2 = 3x^2 \]
となる。
\( 3x^2 \) が、平方数でなければゲームオーバになるので、\( x^2 \) は3の倍数でなければならないが、このルールになる枚数を \( x \) は必ずディーラに要求できるので、ゲームは無限に終わらない。

つまり、このルールでは、\( n \) 次元の人数の盤面では、
\[ l^2 = n x^2 \]
となるので、\( n \) が平方数であれば、ゲームは無限に終わらない。

ここまではごく平凡でつまらないゲームだが、一人のプレーヤーをジョーカー \( t \) とする。
かれだけ特殊なルールがあり、\( x \) とは別の任意の数のカードをディラーに要求できる。
\( t \) も自分の手札は平方数でなければならない。
ところが、ディーラは \( t \) の指定した数のカードをディーラの手札からカードディックに戻す。
これが、\( -t \) のマイナスの概念である。
この場合でも、ディックの枚数は平方数でなければならない。
平方数にできなければ、ゲームオーバーとなる。

このルールはプレーヤ＋ジョーカーの人数 \( n \) が特定の人数の時にのみゲームオーバーが回避できる。
\[ l^2 = (n - 1) x^2 - t^2 \]
が成立する \( n \) である。
これは、
\[ n = 1, 2, 3, 5, 6, 9, 10, 11, 14, 17, 18, 19, 21, 26, 27, 30, 33, 35, 37, \dots \]
となる。

\section{離散ミンコフスキー計量と対称性の仮定 (Discrete Minkowski Metric)}
\(D\) 次元時空における平坦なミンコフスキー計量の線素 \(ds^2\) を考える。光速 \(c=1\) とし、時間軸を1次元、空間軸を \(N\) 次元（\(D = N+1\)）とする。
\[ ds^2 = \sum_{i=1}^{N} (dx_i)^2 - dt^2 \]
時空がプランク長のような最小単位の整数倍で構成される離散格子であると仮定し、すべての変位 \(dx_i, dt\) および不変量 \(ds\) が正の整数をとると制約する。ここで、\(ds\) を \(l\) と表記する。

空間の等方性（isotropy）を維持するため、特定の慣性系において空間の各次元への変位量が等しい対称的な測地線を考える。すなわち、任意の \(i\) について \(dx_i = x \in \mathbb{Z}^+\) とする。このとき、方程式は以下の形に簡約化される。
\[ l^2 = N x^2 - t^2 \implies l^2 + t^2 = N x^2 \]

\section{次元 \(N\) に対するディオファントス的制約 (Diophantine Constraints on Dimension)}
前節で導出された方程式が非自明な解 \((l, t, x) \in (\mathbb{Z}^+)^3\) を持つための、空間次元 \(N\) の条件を特定する。

\begin{theorem}
方程式 \(l^2 + t^2 = N x^2\) が正の整数の解を持つための必要十分条件は、\(N\) の素因数分解において、\(p \equiv 3 \pmod 4\) を満たす任意の素数 \(p\) の指数が偶数であることである。
\end{theorem}

\begin{proof}
背理法およびフェルマーの無限降下法（Method of infinite descent）を用いる。

まず、\(N\) が素因数 \(p \equiv 3 \pmod 4\) を奇数乗で含むと仮定し、\(N = p^{2k+1}M\)（ただし \(p\) と \(M\) は互いに素、\(k \ge 0\)）と置く。このとき、方程式は以下のように表される。
\[ l^2 + t^2 = p^{2k+1} M x^2 \]

フェルマーの二平方和定理より、ある素数 \(p \equiv 3 \pmod 4\) が二つの平方数の和 \(l^2 + t^2\) を割り切る場合、\(p\) は \(l\) と \(t\) の両方を割り切らなければならない。
したがって、\(l = p l'\)、\(t = p t'\) （\(l', t' \in \mathbb{Z}^+\)）と置換できる。これを元の方程式に代入すると、
\[ p^2(l'^2 + t'^2) = p^{2k+1} M x^2 \]
両辺を \(p^2\) で割ると、
\[ l'^2 + t'^2 = p^{2k-1} M x^2 \]
となる。

この操作を繰り返す（\(k\) 回繰り返す）ことで、最終的に右辺の \(p\) の指数は1となり、以下の式を得る。
\[ L^2 + T^2 = p M x^2 \]
ふたたび右辺は \(p\) の倍数であるため、左辺の \(L, T\) も \(p\) の倍数でなければならない。\(L = p L'\)、\(T = p T'\) と置換すると、
\[ p^2(L'^2 + T'^2) = p M x^2 \implies p(L'^2 + T'^2) = M x^2 \]
ここで、左辺は \(p\) の倍数であるため、右辺の \(M x^2\) も \(p\) の倍数でなければならない。しかし \(M\) は \(p\) と互いに素（\(p \nmid M\)）であるため、素数 \(p\) は \(x^2\) を割り切り、結果として \(x\) も \(p\) を因数に持つ（\(x = p x'\)）。

これを上式に代入すると、
\[ p(L'^2 + T'^2) = M p^2 x'^2 \implies L'^2 + T'^2 = M p x'^2 \]
これは \(N = p M\) なので、
\[ L'^2 + T'^2 = N x'^2 \]
となる。

これは全く同じ形式の方程式であるが、解全体のスケールは明らかに小さくなっている（\(L' < l\)、\(T' < t\)、\(x' < x\)）。
正の整数解 \((l, t, x)\) が存在するならば、この操作により無限により小さな正の整数解の組を作り出せることになるが、これは自然数の整列性（無限降下法）に矛盾する。
したがって、正の整数の解は存在せず、自明解 \((0, 0, 0)\) のみに限られる。
\end{proof}

\section{物理的解釈と特定の次元での挙動 (Physical Interpretation)}
定理 3.1 を適用すると、等方的な離散測地線が存在するための空間次元 \(N\) に強い制約が生じることがわかる。

\begin{itemize}
    \item \textbf{\(N=3\)（\(D=4\)）}：\(3 \equiv 3 \pmod{4}\) であるため、定理の条件を満たさず、非自明な正整数解は存在しない。
    \item \textbf{\(N=10\)（\(D=11\)）}：解が存在する（例：\(l=1, t=3, x=1\)）。
    \item \textbf{\(N=25\)（\(D=26\)）}：解が存在する（例：\(l=3, t=4, x=1\)）。
\end{itemize}

これらの結果は、離散格子モデルにおける算術的制約が時空次元の選択に影響を与える可能性を示唆するが、本研究は特殊相対性理論の平坦極限、すなわち平坦ミンコフスキー時空を整数格子上で離散化した場合に限定して厳密解を与える考察である。

\subsection{質量ゼロの光子（ $l=0$ ）の場合の厳密解}

さらに、質量を持たない光子（固有時 $l=0$）の場合に条件を限定すると、方程式は極めて単純な形に帰着される。
\[ 0 = N x^2 - t^2 \implies t^2 = N x^2 \]
この方程式が正の整数解 $(t, x \in \mathbb{Z}^+)$ を持つための必要十分条件は、$\sqrt{N}$ が整数となること、すなわち空間次元 $N$ が「平方数（1, 4, 9, 16, 25, 36, 49, 64, 81, 100, \dots）」であることに限定される。
この条件下においても、$N=3$ では解が存在しないことがより自明に示されるだけでなく、超弦理論における10次元時空の空間次元 $N=9$ （= $3^2$）、およびボソン弦理論における26次元時空の空間次元 $N=25$ （= $5^2$）が、等方直進する光子が立方格子点を正確に通過できる「完全な平方数」として自然に浮かび上がる。

\section{モデルの位置づけと適用範囲 (Scope and Limitations)}
本研究の中心は「新たな数式プロセスの発見」ではなく、「既知の純粋数学（ガウス整数の性質とフェルマーの二平方和定理）」が、離散ミンコフスキー格子という物理モデルに適用された際に、特定の空間次元に対して課す「時空の数学的制約の発見（気づき）」にある。

本研究は、連続時空や量子重力理論の構築を目的とするものではなく、\textbf{離散ミンコフスキー格子上の単純なディオファントス制約}に限定される。

本稿で用いる「対応関係」や「双対性」という言葉は、数学的構造間の\textbf{部分的・形式的な類似性}を指すに過ぎず、物理的同一性や理論的等価性を主張するものではない。

本研究が主張しないこと：
\begin{itemize}
    \item 新たな数学的公式や解法の発見（基礎は既知のガウス整数環におけるノルムの性質による）
    \item 連続時空における物理理論の構築
    \item 量子化手法、ラグランジアン、ダイナミクスの提示
    \item 弦理論や既存の双対性との直接的な同一性・拡張
\end{itemize}

本研究は、相互作用や曲率を含まない平坦時空の離散版という条件下において、厳密な数論的制約を与えるものとして位置づけられる。

\section{議論}
本稿で扱った単一の制約 \(l^2 + t^2 = N x^2\)（\(K=0\) の場合）は、光錐（質量ゼロの殻）に対応する整数点集合を記述する。この制約上での変換は、離散ローレンツ群などの既存の対称性に帰着される。

より一般に \(l^2 + t^2 - N x^2 = K\)（\(K \in \mathbb{Z}\)）と拡張した場合、異なる \(K\) 間の整数点の対応関係を考察できるが、本稿ではその詳細な分類は行わない。

\section{結論 (Conclusion)}
本稿では、特殊相対性理論の平坦ミンコフスキー時空を整数格子上で離散化した際に現れるディオファントス方程式 \(l^2 + t^2 = N x^2\) を考察し、その非自明正整数解の存在条件をフェルマーの二平方和定理と無限降下法により完全に決定した。

特に4次元時空（\(N=3\)）で解が存在しないという算術的制約が明らかになった一方で、特定の \(N\)（例：10, 25）では解が存在することも示された。

このような数論的制約は、離散時空モデルの根本的な性質として興味深い。本研究は平坦ミンコフスキー時空を整数格子上で離散化した場合に限定して厳密解を与えたものであり、複素平面上のガウス整数のノルムとしての性質（フェルマーの二平方和定理）が時空次元の選択に直接適用できることを見出した。連続理論や量子重力への直接的な応用を意図するものではないが、将来的にこのアプローチをさらに発展させる際の、基礎的かつ厳密な数論的知見として役立てば幸いである。

\end{document}
